\honest{Disputing Definitions}{Eliezer Yudkowsky}{2008}

I have watched more than one conversation turn into a dispute over definitions. The classic example is the riddle: if a tree falls in a forest and no one hears it, does it make a sound? \nb{The riddle was answered by a definition in its earliest known printing. In June 1883 \textsc{The Chautauquan} replied: ``No. Sound is the sensation excited in the ear when the air or other medium is set in motion'' (as quoted by Wikipedia).}

Albert says yes: every tree he has heard fall made a sound, and the world does not change when he is not looking. Barry says no: if no one hears it, how can it be a sound? Barry may have a different intuition, a motive, or a skeptic's habit of stopping at the first counterargument. Albert adds that the energy of the fall goes into heat and vibrations, or the tree would violate conservation of energy; Barry, that without anything with a nervous system able to hear, no one hears a sound. They describe what made their ``sound'' detectors fire or not, but they do not disagree about anything that happens in the forest. Then come insults (``niddlewicking fallumphing pickleplumber''), and neither can back down without losing face.

Next the argument turns to ``by definition.'' What is true by definition is true in all possible worlds, so it cannot tell you which world you live in. Albert argues that a microphone records a ``sound file'' of vibrations, not of brain activity; that feels like support, but for the word's meaning, a different question from what happens in the forest, and the shift passes unnoticed. Then Barry reaches for the dictionary. I could go into the forest, derive the wave equation, or study the ear and the auditory cortex. Why a dictionary? ``Are the editors of the dictionary expert botanists, expert physicists, expert neuroscientists?'' Merriam-Webster gives both senses, and the two cry, ``Consarned dictionary!'' Dictionary editors ``are historians of usage, not legislators of language''; they record what people seem to mean, and more than one usage gets more than one definition. Albert appeals to common usage; Barry replies that he can define a word as he likes if he is consistent, that his meaning is in the dictionary, and who made Albert the judge of common usage? \nb{S.~I. Hayakawa wrote in 1949: ``The writer of a dictionary is a historian, not a lawgiver.''}

Albert and Barry agree on virtually everything in the forest, and it brings no feeling of agreement. Arguing about definitions is a garden path, which no one would take if they saw where it led. Each thinks he is defending the standard definition against the other's attempt to sneak in his own.

So I travel fifteen minutes into the past and ask them, before the argument starts, whether ``sound'' should mean vibrations, experience, or both. Albert says it does not matter so long as you are consistent; Barry says to flip a coin. Then I give them a remedy. Describe the event in lower-level terms, acoustic vibrations or auditory experiences. Or give each meaning a new word, ``alberzle'' and ``bargulum,'' so that neither side loses face. And keep track of some testable proposition. Barry then answers the riddle: ``It makes an alberzle but not a bargulum.'' \nb{William James told the same story in \textsc{Pragmatism} (1907), about a man circling a tree while a squirrel keeps the trunk between them. He separated two senses of ``going round'' and wrote: ``Make the distinction, and there is no occasion for any farther dispute.'' David Chalmers later called the renaming step the ``subscript gambit'' (2011).}

This remedy does not settle every dispute over categories, ``But it destroys a substantial fraction.'' I give no count. \nb{The next day's post, ``The Argument from Common Usage,'' names the cases it does not settle.}
